%% RNAAS manuscript for UHI v5.2-final
%% AASTeX v7.0.2 (aastex702.cls, 2025/05/09)
%% Compiled class files are vendored in this directory.
\documentclass[linenumbers]{aastex702}

\shorttitle{Unified Habitability Index v5.2}
\shortauthors{Viridis}

%% Do not guess dates; leave submission dates unset.
\submitjournal{RNAAS}

\begin{document}

\title{A Probabilistic Unified Habitability Index Coupling Planetary Viability with JWST and ELT/PCS Observability across 6,112 Exoplanets}

\author[0009-0004-5620-5352,gname=Effy,sname=Viridis]{Effy Viridis}
\affiliation{Independent Researcher, El Salvador}
\email[show]{effywhitedva@gmail.com}
\correspondingauthor{Effy Viridis}

\begin{abstract}
Traditional habitability metrics such as the Earth Similarity Index (ESI) rank exoplanets by bulk physical analogy without accounting for atmospheric characterization feasibility. We present version 5.2-final of the Unified Habitability Index (UHI), a probabilistic, hurdle-based framework evaluated via 10,000 Monte Carlo draws per planet across 6,112 scored exoplanets from the NASA Exoplanet Archive. UHI combines seven physical, climatic, stellar, and instrumental sub-indices with Dirichlet weight propagation, synchronous M-dwarf climate parameterizations, a hybrid runaway greenhouse veto, and physical observability models for JWST and ELT/PCS. At a $5\sigma$ detection threshold, 515 planets remain viable ($P_{\rm alive} > 0$). Synthetic Solar System anchors at 10\,pc pass cleanly (Earth at the 99.30th (non-transiting) and 99.26th (transiting) percentiles; Venus and Mars $= 0$). Compared against ESI, UHI filters out 71.2\% of high-ESI ($\geq 0.8$) targets that are currently unobservable while elevating seven nearby non-transiting worlds into the Top 20 via reflected-light imaging.
\end{abstract}

%% ``Monte Carlo methods'' has no exact computational UAT match
%% (UAT 2238 is a laboratory reaction-model concept). Closest: Computational methods (1965).
\keywords{\uat{Astrobiology}{74} --- \uat{Computational methods}{1965} --- \uat{Direct imaging}{387} --- \uat{Exoplanets}{498} --- \uat{Habitable planets}{695} --- \uat{Transmission spectroscopy}{2133}}

\section{Introduction}

Prioritizing exoplanets for atmospheric characterization requires balancing intrinsic planetary habitability against realistic instrumental detectability. Widely used similarity indices such as the Earth Similarity Index \citep[ESI;][]{SchulzeMakuch2011} and classical Habitable Zone (HZ) boundaries \citep{Kopparapu2013,Kopparapu2014} successfully identify Earth-sized temperate candidates, yet they treat characterization feasibility as an afterthought. Consequently, many top-ranked analogs lie beyond the reach of current or near-future observatories. Here we summarize version 5.2-final of the Unified Habitability Index (UHI), a fully open, reproducible Monte Carlo framework designed to bridge intrinsic habitability and observational priority for both transiting (JWST) and non-transiting (ELT/PCS) exoplanets.

\section{Methodology}

We evaluate 6,112 scored exoplanets from a dated snapshot (2026-11-18) of the NASA Exoplanet Archive Planetary Systems Composite Parameters table \citep[\texttt{pscomppars}; 6,470 rows, 358 excluded for missing core parameters;][]{Akeson2013} using 10,000 Monte Carlo draws per planet (seed 42). Measurement uncertainties are sampled from split-normal distributions, with missing masses or radii propagated via the probabilistic mass--radius relation of \citet{Chen2017}.

The composite index is defined as a weighted geometric mean, ${\rm UHI} = \prod S_i^{w_i}$, over seven sub-indices---physical ($S_{\rm fis}$), climate ($S_{\rm clima}$), atmospheric retention ($S_{\rm atm}$), stellar environment ($S_{\rm estrella}$), geophysical ($S_{\rm geo}$), magnetic ($S_{\rm mag}$), and observability ($S_{\rm obs}$)---with weights drawn per realization from ${\rm Dirichlet}(\alpha = [2, 2, 2, 2, 1, 1, 2])$. A numerical floor $\varepsilon = 10^{-12}$ precedes log-aggregation, imposing an exact veto (${\rm UHI} = 0$) whenever any $S_i \leq \varepsilon$. Key architectural features of v5.2-final include:

\begin{itemize}
\item \textbf{Orthogonal Atmospheric--Observational Decoupling:} $S_{\rm atm}$ models intrinsic atmospheric retention via the bolometric Cosmic Shoreline \citep{Zahnle2017}, while all instrumental detection terms reside strictly in $S_{\rm obs} = P_{\rm obs} \times {\rm Availability(tier)}$. Transiting observability combines JWST transmission (TSM-scaled with a 20\,ppm NIRISS systematic floor; \citealt{Kempton2018,Greene2016}) and secondary eclipse emission (ESM), whereas non-transiting targets use a residual-halo photon-noise limited ELT/PCS reflected-light model \citep{Kasper2021} with inner working angle ${\rm IWA} \sim U(2, 4)\,\lambda/D$ at 0.75\,$\mu$m. The primary threshold is $5\sigma$, with a $3\sigma$ comparison variant.

\item \textbf{Synchronous M-Dwarf Climate and Hybrid Runaway Veto:} $S_{\rm clima}$ couples a 1-D two-mode energy balance model \citep{North1981} and $T_{\rm eff}$-dependent Bond albedos \citep{Shields2013} with 3-D GCM-calibrated parameterizations for tidally locked M-dwarf planets ($T_{\rm eff} < 3800$\,K): substellar cloud albedos \citep{Yang2013} smoothly blended via a linear ramp $w_{\rm eff}(S_0) = {\rm clip}((S_0 - 800)/200,\, 0,\, 1)$, inner-edge expansion $k_{\rm in} \sim U(1.5, 2.0)$ \citep{Kopparapu2016}, and an outer-edge ${\rm CO}_2$ surface-pressure floor \citep{Turbet2018}. For non-synchronous or $T_{\rm eff} \geq 3800$\,K planets, a hybrid runaway greenhouse veto enforces $f_{\rm HZ} = 0$ whenever absorbed flux exceeds 282\,W\,m$^{-2}$ \citep{Goldblatt2013,Goldblatt2014}.

\item \textbf{Smooth Mean Molecular Weight and Hurdle Ranking:} Across the radius valley \citep{Fulton2017}, $\mu(R_p)$ transitions log-sigmoidally between secondary (28 or 44\,g\,mol$^{-1}$) and ${\rm H}_2$/He (2.3\,g\,mol$^{-1}$) envelopes centered at $R_0 \sim N(1.65, 0.08)\,R_\oplus$ ($\Delta R = 0.12\,R_\oplus$), overridden to secondary atmospheres when measured densities exceed 5.5\,g\,cm$^{-3}$ ($M/\sigma_M > 3$). Because multiplicative vetoes yield bimodal posteriors, planets are ranked by a continuous hurdle metric: ${\rm UHI}_{\rm rank} = P_{\rm alive} \times {\rm median}({\rm UHI}\mid {\rm UHI} > 0)$, accompanied by $P({\rm top\,10})$ and 90\% rank credible intervals.
\end{itemize}

\section{Results, Validation, and Benchmarking}

Across the 6,112 planets, 515 remain viable ($P_{\rm alive} > 0$) at both $5\sigma$ and $3\sigma$ (Spearman $\rho = 0.985$ between variants). Synthetic Solar System twins injected at 10\,pc validate the pipeline: Earth ranks at the 99.30th percentile (\#44, non-transiting) and 99.26th (\#46, transiting) with $P_{\rm alive} = 1.00$, whereas Venus and Mars yield ${\rm UHI} = 0$ ($P_{\rm alive} = 0$).

Table~\ref{tab:top10} summarizes the Top 10 targets at $5\sigma$ alongside their $3\sigma$ rankings and corrected ESI values \citep{SchulzeMakuch2011}. At $5\sigma$, temperate sub-Neptunes around M dwarfs---TOI-4336 A b (${\rm UHI}_{\rm rank} = 0.562$) and TOI-7166 b (0.547)---lead due to large ${\rm H}_2$-rich scale heights and an upper radius trapezoid limit of $f_{R,d} = 2.5\,R_\oplus$. Relaxing to $3\sigma$ (a proxy for deeper JWST campaigns) elevates rocky worlds: TRAPPIST-1 d climbs from \#5 (0.436) to \#1 (0.609), and TRAPPIST-1 e rises from \#17 (0.272) to \#9 (0.428). This shift directly reflects the observational trade-off identified by a 51-parameter Sobol sensitivity analysis \citep[SALib;][]{Herman2017} performed after v5 and not rerun for v5.1 or v5.2. The quoted total-effect indices are means across that run's 13 per-planet outputs and 3 aggregate metrics: the detection threshold $\mathtt{snr\_thr}$ (mean $S_T = 0.45$) and $f_{R,d}$ (mean $S_T = 0.15$) dominate variance, while the newly added transition parameters ($R_0$, $\Delta R$, $k_{\rm in}$) have mean $S_T < 0.01$. The Sobol analysis does not cover the runaway-greenhouse flux threshold $F_{\rm rg}$ or the substellar-cloud mixing ramp $w_{\rm eff}$.

\begin{deluxetable}{cclcccccc}
\tablecaption{Top 10 UHI v5.2-final target exoplanets at $5\sigma$ detection threshold with $3\sigma$ comparison metrics, conditional posterior distributions, and classical ESI.\label{tab:top10}}
\tablewidth{0pt}
\tabletypesize{\scriptsize}
\tablehead{
\colhead{Rank ($5\sigma$)} & \colhead{Rank ($3\sigma$)} & \colhead{Planet} &
\colhead{UHI$_{\rm rank}$ ($5\sigma$)} & \colhead{$P_{\rm alive}$} &
\colhead{Cond.\ median [IQR]} & \colhead{$P$(top 10)} & \colhead{Tier} & \colhead{ESI$_{\rm SM11}$}
}
\startdata
1 & 2 & TOI-4336 A b & 0.562 & 0.99 & 0.567 [0.460--0.665] & 0.60 & T0 & 0.741 \\
2 & 3 & TOI-7166 b & 0.547 & 0.99 & 0.554 [0.456--0.647] & 0.56 & T1 & 0.816 \\
3 & 4 & GJ 273 b & 0.477 & 0.87 & 0.546 [0.424--0.642] & 0.47 & T2 & 0.847 \\
4 & 6 & K2-332 b & 0.449 & 0.91 & 0.494 [0.393--0.584] & 0.36 & T2 & 0.786 \\
5 & 1 & TRAPPIST-1 d & 0.436 & 1.00 & 0.436 [0.261--0.602] & 0.35 & T0 & 0.905 \\
6 & 8 & Wolf 1061 c & 0.434 & 0.82 & 0.529 [0.426--0.625] & 0.41 & T2 & 0.823 \\
7 & 11 & Proxima Cen d & 0.389 & 0.84 & 0.463 [0.356--0.566] & 0.28 & T2 & 0.876 \\
8 & 5 & TOI-2094 b & 0.377 & 0.99 & 0.382 [0.229--0.521] & 0.23 & T2 & 0.813 \\
9 & 17 & GJ 1061 d & 0.342 & 0.64 & 0.537 [0.429--0.611] & 0.33 & T2 & 0.884 \\
10 & 15 & K2-125 b & 0.333 & 0.84 & 0.396 [0.271--0.502] & 0.17 & T2 & 0.767 \\
\enddata
\end{deluxetable}

The 90\% rank credible intervals are wide (e.g., 1--30 for the leading target) and are provided in the archive.

Benchmarking against ESI$_{\rm SM11}$ and the Habitable Worlds Catalog \citep[][file last updated 2024-03-12]{Mendez2024} confirms that UHI provides distinct prioritization (Spearman $\rho = 0.466$ vs.\ ESI$_{\rm SM11}$ and 0.444 vs.\ PHL ESI among viable worlds). Specifically, 71.2\% (47/66) of planets with ESI$_{\rm SM11} \geq 0.8$ have median detection probability $P_{\rm det} < 0.001$ at $5\sigma$ (72.0\%, 18/25 for PHL ESI); UHI demotes these unobservable analogs so that none enter the top 29 (best: Kepler-1653 b at \#30). Simultaneously, UHI rescues 26 nearby non-transiting viable worlds with median $P_{\rm det} \geq 0.5$ via ELT/PCS, placing seven in the Top 20 (GJ 273 b, Wolf 1061 c, Proxima Cen d, GJ 1061 d, Ross 128 b, Proxima Cen b, and GJ 433 d).

A 1-D sweep of the runaway veto across $F_{\rm rg} \in [275, 295]$\,W\,m$^{-2}$ preserves Venus $P_{\rm alive} = 0$, an identical Top 10, and Spearman $\rho \geq 0.996$ among viable worlds. Sweeping the thermal boundary to $T_{\rm rg} = 3700$--$3750$\,K swaps \#10 K2-125 b ($T_{\rm eff} = 3654$\,K, ${\rm UHI}_{\rm rank} = 0.333$) with \#11 Ross 128 b (0.323) while leaving the other nine unchanged; we retain $T_{\rm rg} = 3800$\,K without post-hoc tuning and report this boundary sensitivity transparently.

\section{Reproducibility, Software, and Acknowledgments}

All code (MIT license), dated catalog snapshots, machine-readable tables, and figures (CC BY 4.0) for UHI v5.2-final are archived at \url{https://github.com/effywhitedva-ops/uhi-habitability-index} and Zenodo (\dataset[doi:10.5281/zenodo.24358422]{https://doi.org/10.5281/zenodo.24358422}). The PHL catalog is fetched at runtime and not redistributed.

% Typographic squeeze so the DOI-records disclosure tweak stays on 3 pages.
\vspace{-14pt}
\begin{acknowledgments}
This research has made use of the NASA Exoplanet Archive, operated by the California Institute of Technology under contract with NASA under the Exoplanet Exploration Program.

\textbf{AI disclosure.} AI agents based on Grok (via Grok Bot; agents ``Shoggothsito'' and ``Aetherios''), with computing tools and under the author's supervision, were used to implement and refactor the Python code, run Monte Carlo and sensitivity simulations on local and Google Colab runtimes, generate tables and figures from the author's data with Matplotlib, audit methods, assist in drafting, editing, and formatting the manuscript, and supplement the literature search. Gemini was used as a reviewer to provide critical feedback on the methodology and manuscript drafts; all suggestions were evaluated by the author. All references were checked against the original sources (DOI records) and {NASA} {ADS}. The author reviewed and approved every change and takes full responsibility for the content. AI tools are not authors.
\end{acknowledgments}

\software{Python 3.13, NumPy, SciPy, pandas, SALib \citep{Herman2017}, Matplotlib, Astropy, PyArrow}

\begin{thebibliography}{99}

\bibitem[Akeson et al.(2013)]{Akeson2013} Akeson, R.~L., Chen, X., Ciardi, D., et al. 2013, PASP, 125, 989

\bibitem[Chen \& Kipping(2017)]{Chen2017} Chen, J., \& Kipping, D. 2017, ApJ, 834, 17

\bibitem[Fulton et al.(2017)]{Fulton2017} Fulton, B.~J., Petigura, E.~A., Howard, A.~W., et al. 2017, AJ, 154, 109

\bibitem[Goldblatt et al.(2013)]{Goldblatt2013} Goldblatt, C., Robinson, T.~D., Zahnle, K.~J., \& Crisp, D. 2013, NatGeo, 6, 661

\bibitem[Goldblatt et al.(2014)]{Goldblatt2014} Goldblatt, C., Robinson, T.~D., Zahnle, K.~J., \& Crisp, D. 2014, NatGeo, 7, 931 (Corrigendum)

\bibitem[Greene et al.(2016)]{Greene2016} Greene, T.~P., Line, M.~R., Montero, C., et al. 2016, ApJ, 817, 17

\bibitem[Herman \& Usher(2017)]{Herman2017} Herman, J., \& Usher, W. 2017, JOSS, 2, 97

\bibitem[Kasper et al.(2021)]{Kasper2021} Kasper, M., Cerpa Urra, N., Pathak, P., et al. 2021, Msngr, 182, 38

\bibitem[Kempton et al.(2018)]{Kempton2018} Kempton, E.~M.-R., Bean, J.~L., Louie, D.~R., et al. 2018, PASP, 130, 114401

\bibitem[Kopparapu et al.(2013)]{Kopparapu2013} Kopparapu, R.~K., Ramirez, R., Kasting, J.~F., et al. 2013, ApJ, 765, 131

\bibitem[Kopparapu et al.(2014)]{Kopparapu2014} Kopparapu, R.~K., Ramirez, R.~R., SchottelKotte, J., et al. 2014, ApJ, 787, L29

\bibitem[Kopparapu et al.(2016)]{Kopparapu2016} Kopparapu, R.~K., Wolf, E.~T., Haqq-Misra, J., et al. 2016, ApJ, 819, 84

\bibitem[{PHL @ UPR Arecibo}(2024)]{Mendez2024} PHL @ UPR Arecibo 2024, The Habitable Worlds Catalog (file last updated 2024-03-12), \url{https://phl.upr.edu/hwc}

\bibitem[North et al.(1981)]{North1981} North, G.~R., Cahalan, R.~F., \& Coakley, J.~A. 1981, RGeo, 19, 91

\bibitem[Schulze-Makuch et al.(2011)]{SchulzeMakuch2011} Schulze-Makuch, D., M\'endez, A., Fair\'en, A.~G., et al. 2011, Astrobiology, 11, 1041

\bibitem[Shields et al.(2013)]{Shields2013} Shields, A.~L., Meadows, V.~S., Bitz, C.~M., et al. 2013, Astrobiology, 13, 715

\bibitem[Turbet et al.(2018)]{Turbet2018} Turbet, M., Bolmont, E., Leconte, J., et al. 2018, A\&A, 612, A86

\bibitem[Yang et al.(2013)]{Yang2013} Yang, J., Cowan, N.~B., \& Abbot, D.~S. 2013, ApJL, 771, L45

\bibitem[Zahnle \& Catling(2017)]{Zahnle2017} Zahnle, K.~J., \& Catling, D.~C. 2017, ApJ, 843, 122

\end{thebibliography}

\end{document}
